\documentclass[10pt,a4paper]{amsart}
\usepackage[margin=19mm]{geometry}
\usepackage{amsmath,amssymb,mathtools}
\usepackage[scr=rsfs,cal=euler]{mathalfa}
\usepackage{xfrac}
\usepackage[unicode,hidelinks,bookmarks=false]{hyperref}
\newcommand{\ie}{i.e.\ }
\newcommand{\cf}{cf.\ }
\newcommand{\resp}{respectively\ }
\newcommand{\Subsection}{\subsection}
\providecommand{\nobreakdash}{\nobreak\hbox{-}}
\setlength{\emergencystretch}{2em}
\multlinegap0pt
\usepackage{bm}
\usepackage{bbm}
\usepackage{dsfont}
\usepackage{xspace}
\usepackage{cancel}
\usepackage{mathtools}
\usepackage{amsthm}
\makeatletter
\@ifundefined{definition}{\newtheorem{definition}{Definition}}{}
\@ifundefined{remark}{\newtheorem{remark}[definition]{Remark}}{}
\@ifundefined{theorem}{\newtheorem{theorem}[definition]{Theorem}}{}
\makeatother
\title[A Conservation Law for Equilibrium Propagation and Coupled L]{A Conservation Law for Equilibrium Propagation and Coupled Learning}
\author{Joshua A. McGinnis, Adam G. Kline, Yoichiro Mori}
\date{Source: arXiv:2606.15444v1 (CC BY 4.0)}
\begin{document}
\maketitle

\noindent\emph{Source and licence.} A Conservation Law for Equilibrium Propagation and Coupled Learning. Version arXiv:2606.15444v1 (CC BY 4.0), distributed under the Creative Commons Attribution 4.0 International licence, as recorded in the arXiv licence metadata for this exact version and shown on the abstract page https://arxiv.org/abs/2606.15444v1. Full rights notice: licenses/arxiv-2606.15444-CC-BY-4.0.txt.\par\smallskip

\noindent\textit{Editorial scope.} Counted unit: the Theorem whose source label is \texttt{thm:crossbar} in the article (source0.tex lines 1185--1190, source label \texttt{thm:crossbar}), delivered together with the article's own complete original proof of it (source0.tex lines 1192--1279, one \texttt{proof} environment of 88 lines). No mathematics is written, repaired or re-derived: statement and proof are the article's own text line for line. Required context retained uncounted: Thm:Cond (lines 354--366); Rmk:EP\_Crossbar (lines 1281--1290). Every definition, display and label of those lines is retained unchanged. Omitted from this package and not redistributed: 1--353, 367--1184, 1191--1191, 1280--1280, 1291--1781. Nothing inside a retained excerpt is omitted or altered: every retained body is delivered exactly as the article's lines are written, so no omission receipt of any kind accompanies this package. Citations: the retained excerpts carry no citation command at all, so no bibliography entry of the article is transcribed and no reference is redistributed. Reference adaptation: none was needed and none was made; every reference inside the delivered unit resolves inside it, so no unresolved reference or citation is produced. Printed slips kept literal: no printed word, symbol, hypothesis or number of the counted statement or of its proof was altered. Layout and class: the article's own class is \texttt{elsarticle} with packages including amsmath, amssymb, amsthm, tikz, circuitikz, mathrsfs, bbm, mathtools, enumitem, hyperref, xcolor, caption, ulem, lineno; the wrapper is the lane's pinned \texttt{amsart} editorial wrapper at 10pt on A4, which restates the theorem-like environments the retained text needs and adds the article's own macro definitions that the retained text needs (none), with extra wrapper packages (bm, bbm, dsfont, xspace, cancel, mathtools). The delivered numbering is the unit's own. Assets: no retained span contains a figure environment, a graphics-inclusion command, a PDF-page inclusion, a table or a TikZ picture; no image file is copied into this package, the package declares no asset dependency and no image file is redistributed with it. Licence: version arXiv:2606.15444v1 of this article is distributed under the Creative Commons Attribution 4.0 International licence, as recorded in the arXiv licence metadata for this exact version and shown on the abstract page https://arxiv.org/abs/2606.15444v1. A full rights notice is delivered at licenses/arxiv-2606.15444-CC-BY-4.0.txt. Characters: every retained excerpt is pure ASCII, so the delivered engine, XeLaTeX with its default Latin Modern text font, reports no missing character.
\begin{theorem}[Conductance ``mass'' is conserved]
    Consider training the conductances in a linear resistive circuit.  
    In this case the energy takes the form
    \[
        G(x;\kappa) = \sum_{i=1}^{M} \kappa_i\,\widetilde{G}_i(x),
    \]
    where $\kappa_i$ represents the $i$th conductance.
    Then the conserved quantity is
    \[
        K(\kappa) = \frac{1}{2}\sum_{i=1}^{M} \kappa_i^2.
    \]
    \label{Thm:Cond}
\end{theorem}

\begin{remark}
The EP dynamics for the crossbar are
$\dot{\kappa}_i = -r\,\bar{\kappa}^{-1}(x_i - \mu)$,
which differ from the CL dynamics only by the factor $\bar{\kappa}^{-1}$.
The loss dynamics become
$\dot{\Phi} = -2\Phi\,\bar{\kappa}^{-2}\,V_{\mathcal{A}}$,
and the same argument gives exponential convergence with rate
$\lambda_{\mathrm{EP}} = 2\,V_{\min}/(2NK) = V_{\min}/(NK)$.
\label{Rmk:EP_Crossbar}
\end{remark}

\begin{theorem}\label{thm:crossbar}                   
  In the setting described above, for both CL and EP, the loss                 
  \[\Phi(t) = \dfrac{1}{2}\big(\mu(t) - w\big)^2 \]
  converges to zero exponentially. This holds even if some conductances reach
  zero and are removed during training.
  \end{theorem}

  \begin{proof}

  We give the proof for CL; the EP case is addressed in Remark~\ref{Rmk:EP_Crossbar} below. Computing the CL update gives
\[
\dot{\kappa}_i
= \lim_{\eta \to 0}
  \dfrac{1}{\eta}\Big(\partial_{\kappa_i}G\big(\mu + \eta r;\, \kappa\big)
  - \partial_{\kappa_i}G\big(\mu;\, \kappa\big)\Big)
= -r\,(x_i - \mu).
\]

Write $\mathcal{A}(t) := \{i : \kappa_i(t) > 0\}$ for the set of
\emph{active} edges.  Since
$\partial_{\kappa_i}\mu = \bar{\kappa}^{-1}(x_i - \mu)$, the loss dynamics
restricted to active edges are
\[
\dot{\mu}
= \bar{\kappa}^{-1}\sum_{i \in \mathcal{A}} \dot{\kappa}_i\,(x_i - \mu)
= -\dfrac{r}{\bar{\kappa}} \sum_{i \in \mathcal{A}} (x_i - \mu)^2.
\]
Writing $V_{\mathcal{A}} := \sum_{i \in \mathcal{A}} (x_i - \mu)^2$, we obtain
\begin{equation}\label{eq:crossbar_loss}
\dot{\Phi} = r\,\dot{\mu} = -\dfrac{2\Phi}{\bar{\kappa}}\, V_{\mathcal{A}}
\leq 0.
\end{equation}
We need to show that $V_{\mathcal{A}}$ remains bounded below.

Whenever an edge's conductance reaches zero, we remove it from the circuit.
Since $\kappa_i = 0$ at the moment of removal, neither $\mu$ nor $\bar{\kappa}$
changes discontinuously, and the removed edge contributes nothing to the
conserved quantity $K = \frac{1}{2}\sum \kappa_i^2$.
The ODE continues on the reduced system.

If $r = 0$ then $\Phi = 0$ and the target is already achieved.
Otherwise, since $\dot{\Phi} \leq 0$, the error $r$ cannot change sign.
We treat the case $r < 0$ (i.e.\ $\mu(t) < w$ for all $t$);
the case $r > 0$ is symmetric. Since $r < 0$, the dynamics $\dot{\kappa}_i = -r\,(x_i - \mu)$ give:
\begin{itemize}
\item $x_i > \mu(t)$:\; $\dot{\kappa}_i > 0$ (growing),
\item $x_i < \mu(t)$:\; $\dot{\kappa}_i < 0$ (shrinking, may hit $0$).
\end{itemize}

Since $\mu(t) < w < x_1$, every edge with $x_i \geq w$ satisfies
$x_i > \mu(t)$, so $\dot{\kappa}_i > 0$.
These edges are growing and never reach zero. Moreover, at least one active edge below $w$ exists, or else $\mu \geq w$, contradicting $r < 0$.
So at least one active edge has $x_i < w$, and it is clearly distinct
from $x_1$.

Among active edges, edge~$1$ (with input $x_1$) satisfies $x_1 > \mu$, so
\[
V_{\mathcal{A}} \geq (x_1 - \mu)^2 \geq (x_1 - w)^2 > 0,
\]
since $\mu < w < x_1$.  This bound is uniform in time.
In the case $r > 0$, the symmetric argument shows that $x_N$ always
survives and gives $V_{\mathcal{A}} \geq (x_N - w)^2 > 0$.  In either case,
\[
V_{\min} := \min\!\big((x_1 - w)^2,\, (x_N - w)^2\big) > 0.
\]

By Theorem \ref{Thm:Cond} 
$\|\kappa\|^2 = \sum \kappa_i^2 = 2K$ is conserved.
Since all $\kappa_i \geq 0$,
\[
\bar{\kappa}^2 = \Big(\sum_i \kappa_i\Big)^2
\geq \sum_i \kappa_i^2 = 2K,
\]
and by Cauchy--Schwarz,
\[
\bar{\kappa} \leq \sqrt{N}\,\|\kappa\| = \sqrt{2NK},
\]
so $\bar{\kappa}$ remains bounded above and below for all time.

Combining \eqref{eq:crossbar_loss} with the bounds on $\bar{\kappa}$ and
$V_{\mathcal{A}}$:
\begin{equation}
\label{eq:Phi_dot_crossbar}
\dot{\Phi}
= -\dfrac{2\Phi}{\bar{\kappa}}\, V_{\mathcal{A}}
\leq -\dfrac{2\,V_{\min}}{\bar{\kappa}}\,\Phi
\leq -\dfrac{2\,V_{\min}}{\sqrt{2NK}}\,\Phi
=: -\lambda\,\Phi.
\end{equation}
By Gronwall's inequality,
\[
\Phi(t) \leq \Phi(0)\, e^{-\lambda t}.
\]
The loss converges to zero exponentially, and hence $\mu(t) \to w$.
\end{proof}

\end{document}
